\title{xLSTM: Extended Long Short-Term Memory}

\begin{abstract}
The foundational mechanisms behind Long Short-Term Memory (LSTM)—specifically, the constant error carousel and gating—were formulated during the 1990s, establishing LSTM as a cornerstone in deep learning and paving the way for the initial Large Language Models (LLMs). Despite these advancements, the emergence of Transformer architectures, distinguished by their scalable and parallelizable self-attention mechanisms, has led to a paradigm shift, placing LSTMs at a disadvantage at large scale. This work revisits a fundamental inquiry: To what extent can language modeling be advanced by scaling LSTMs to billions of parameters, making use of contemporary techniques developed for modern LLMs, while also addressing established weaknesses inherent to LSTMs? To this end, we propose exponential gating alongside tailored normalization and stabilization strategies. Furthermore, we reconfigure the memory system of the LSTM, producing: (i) sLSTM, characterized by a scalar memory, a scalar update procedure, and an updated mixing mechanism for memory, and (ii) mLSTM, which is inherently parallelizable due to its matrix memory and a rule for updating covariance. The integration of these innovations as residual block modules creates xLSTM blocks, which are in turn used to build deep xLSTM network architectures via residual stacking. The application of exponential gating and redesigned memory architectures enables xLSTM networks to compete effectively with state-of-the-art Transformer and State Space Models, both with regard to performance and scalability.\\
Code available at: \url{https://github.com/NX-AI/xlstm}
\end{abstract}

\section{Introduction}
\label{sec:intro}

The Long Short-Term Memory (LSTM) ideas~\citep{Hochreiter:91,Hochreiter:97nips,Hochreiter:97}, i.e., the constant error carousel and gating, were introduced to overcome the vanishing gradient problem of recurrent neural networks~\citep{Hochreiter:91,Hochreiter:00book}:

\begin{align}
\label{eq:lstm_idea}
  c_t \ = \ f_t \ c_{t-1} \ + \ 
          i_t \  z_t \ , \quad  
          h_t \ = \ o_t \ \psi({c_t}) \ . 
\end{align}

The constant error carousel refers to the process by which the cell state $c_{t-1}$ (depicted in green) is incrementally updated through the addition of the cell input $z_t$, with this process regulated by sigmoid gating mechanisms (shown in blue). Specifically, the input gate $i_t$ and the forget gate $f_t$ determine how this additive modification occurs, while the output gate $o_t$ governs how information is transmitted from the memory cell, resulting in the hidden state $h_t$. Prior to producing the hidden state, the cell state undergoes normalization or compression via the function $\psi$, after which the output gate selects the information to form the final hidden state.

LSTMs have found widespread success across a broad spectrum of domains \citep{Hochreiter:01,Hochreiter:07,Schmidhuber:15} and remained dominant in text generation tasks prior to the introduction of Transformers in 2017~\citep{Vaswani:17}. Their utility has been validated through a variety of sequence-centric applications, including but not limited to text generation~\citep{Graves:13,Karpathy:15blog}, handwriting synthesis~\citep{Graves:13}, sequence-to-sequence translation~\citep{Sutskever:14nips}, computer program evaluation~\citep{Zaremba:14arxiv}, image caption generation~\citep{Karpathy:15,Hossain:19}, source code production~\citep{Karpathy:15blog}, modeling of rainfall-runoff dynamics~\citep{Kratzert:18,Kratzert:19}, and hydrologic models used for flood warning systems~\citep{Nearing:24}. Within reinforcement learning, LSTMs maintain a position as leading sequence models, underpinning notable systems such as AlphaStar for StarCraft II~\citep{Vinyals:17short}, OpenAI Five for Dota 2~\citep{OpenAI:19}, and controllers for magnetic confinement in nuclear fusion~\citep{Degrave:22short}. Their proficiency lies in the ability to derive abstract representations, skillfully capturing and preserving semantic content within their memory cells~\citep{Karpathy:15blog}. This capability is illustrated by the presence of neurons sensitive to numbers and syntax~\citep{Lakretz:19}, linguistic features~\citep{Bau:19}, and sentiment~\citep{Radford:17arxiv}. Even today, LSTMs continue to play a vital role in critical real-world applications~\citep{Degrave:22short,Nearing:24}, evidencing their enduring relevance.

Although LSTMs have achieved remarkable success, they suffer from three significant drawbacks: (i) Lack of ability to revisit and alter prior storage decisions. We illustrate this deficiency through the {\it Nearest Neighbor Search} task (details provided in Appendix~\ref{sec:appNearestNeighborSearch}): given a reference vector, the model must identify the most similar vector within a sequence and return the associated value at the end of the sequence. As depicted in the left panel of Figure~\ref{fig:lstmProblems}, the mean squared error highlights LSTM's difficulty in updating a previously stored value when a closer match appears later; by contrast, our proposed xLSTM resolves this issue with exponential gating. (ii) Constrained storage capability, since the architecture compresses information into one-dimensional cell states. This constraint is showcased with the {\it Rare Token Prediction} task, presented in the right panel of Figure~\ref{fig:lstmProblems}, which displays the perplexity for predicting tokens at various frequency levels on the Wikitext-103 dataset~\citep{Merity:17}. LSTMs underperform on infrequent tokens as a consequence of their limited memory capacity, whereas xLSTM addresses this problem through matrix-based memory. (iii) Restrictions on parallel computation due to intertwined memory states: the recurrence across timesteps demands strictly sequential processing via hidden-to-hidden state connections.

Such shortcomings have facilitated the ascendance of Transformer models~\citep{Vaswani:17} in the field of language modeling. This raises a natural question: If these limitations were surmounted, and LSTMs were scaled to match the size of contemporary Large Language Models, what levels of performance might we achieve in language modeling?

\section{Extended Long Short-Term Memory}
\label{sec:xLSTM}

To address the shortcomings inherent in traditional LSTM models, the Extended Long Short-Term Memory (xLSTM) architecture implements two significant enhancements over the conventional LSTM scheme detailed in Equation~\eqref{eq:lstm_idea}. These enhancements — specifically, the introduction of exponential gating mechanisms and new memory architectures — contribute two new members to the LSTM family: (i) sLSTM (refer to Section~\ref{sec:sLSTM}), characterized by a scalar memory, scalar memory updates, and the incorporation of memory mixing; and (ii) mLSTM (detailed in Section~\ref{sec:mLSTM}), which employs a memory matrix and updates via covariance (outer product), an operation that facilitates full parallelism. Both variants augment the core LSTM through the use of exponential gating functions. To optimize for parallel computation, mLSTM dispenses with the memory mixing component, thereby removing hidden-hidden recurrent interconnections. Both mLSTM and sLSTM variants support extension to configurations with multiple memory cells. In the sLSTM, memory mixing can occur between individual cells, while additionally supporting multiple heads; in this case, memory mixing is restricted to occur only within cells inside the same head, and not across different heads. This configuration, along with exponential gating, introduces a novel strategy for mixing memory. Conversely, in mLSTM, the constructs of multiple heads and multiple cells are functionally identical.

The assimilation of these LSTM variants into residual block frameworks yields the so-called xLSTM blocks, as described in Section~\ref{sec:xLSTMblock}. Layering these xLSTM blocks in a residual manner gives rise to comprehensive xLSTM architectures, as presented in Section~\ref{sec:xLSTMarchitecure}. For a depiction of the overall xLSTM architecture and its key components, consult Figure~\ref{fig:xlstm_sketch}.

\subsection{Review of the Long Short-Term Memory}
\label{sec:orgLSTM}

The foundational concept of LSTM networks~\citep{Hochreiter:91,Hochreiter:97nips,Hochreiter:97} proposed the use of a scalar memory cell as the main unit for information processing and retention, thereby mitigating the vanishing gradient issue~\citep{Hochreiter:91,Hochreiter:00book} through the mechanism termed the constant error carousel (the cell state update). This memory cell is regulated by three distinct gates: the input gate, the output gate, and the forget gate. The introduction of the forget gate was later provided by \citet{Gers:00}. The memory cell's behavior at time step $t$ is governed by the following update equations:

\begin{align}
c_t \ &= \  f_t \ c_{t-1} \ + \ i_t \ z_t &  & &\text{cell state} \\
h_t  \ &= \ o_t \ \tilde{h}_t \ , & \tilde{h}_t \ &= \ \psi \left( c_t\right)
  &\text{hidden state} \\
z_t \ &= \ \varphi \left( \tilde{z}_t \right) \ , 
  &\tilde{z}_t \ &=  \ \mathbf{w}^\top_{z} \ \mathbf{x}_t \ + \
  r_{z}  h_{t-1} \ + \  b_{z} \ \
  &\text{cell input} \\
i_t \ &= \ \sigma \left( \tilde{i}_t  \right) \ , 
  &\tilde{i}_t \ &= \ \mathbf{w}^\top_{i} \ \mathbf{x}_t \ + \
  r_{i}  \ h_{t-1} \ + \  b_{i} \ \
  &\text{input gate} \\
f_t \ &= \ \sigma \left(  \tilde{f}_t \right) \ , 
  &\tilde{f}_t \ &= \ \mathbf{w}^\top_{f} \ \mathbf{x}_t  \ + \
  r_{f}  \ h_{t-1} \ + \  b_{f} \ \
  &\text{forget gate} \\
o_t \ &= \ \sigma \left( \tilde{o}_t \right) \ , 
  &\tilde{o}_t  \ &= \ \mathbf{w}^\top_{o} \ \mathbf{x}_t \ + \
  r_{o}  \ h_{t-1} \ + \  b_{o} \ \
  &\text{output gate} 
\end{align}

The input weight vectors $\mathbf{w}_{z}$, $\mathbf{w}_{i}$, $\mathbf{w}_{f}$, and $\mathbf{w}_{o}$ are associated, respectively, with the connections from the input $\mathbf{x}_t$ to the cell input, input gate, forget gate, and output gate. In a similar fashion, $r_{z}$, $r_{i}$, $r_{f}$, and $r_{o}$ represent the recurrent weight parameters that connect the previous hidden state $h_{t-1}$ to the corresponding cell input and gates. The terms $b_{z}$, $b_{i}$, $b_{f}$, and $b_{o}$ denote their respective biases. The functions $\varphi$ and $\psi$ define the nonlinearities applied to the cell input and hidden state, most commonly realized as $\tanh$. The nonlinearity $\psi$ in particular is responsible for constraining or rescaling the cell state, thus preventing it from diverging. Sigmoid functions serve as activation functions for all gating mechanisms, with the standard definition given by $\sigma \left( x \right)= 1/(1+\exp(-x))$. Later developments introduced the use of vectorized memory cells $\mathbf{c}_t \in \mathbb{R}^{d}$, where recurrent weight matrices $\mathbf{R} \in \mathbb{R}^{d \times d}$ enable interaction among the outputs from different memory cells~\citep{Greff:15}; additional exposition is available in Appendix~\ref{sec:apporgLSTM}. Ablation analyses confirm the necessity of each element within the memory cell architecture~\citep{Greff:15}.

\subsection{sLSTM}
\label{sec:sLSTM}

To enable LSTMs to modify their memory allocation choices, we incorporate exponential gates (indicated in red), as well as mechanisms for normalization and stabilization. Specifically, both the input and forget gates may employ exponential activation functions. For the purpose of normalization, we define a normalizer state that accumulates the products of the input gate with all subsequent forget gates.

The scalar sLSTM forward pass is:

\begin{align}
\label{eq:slstmforward}
c_t \ &= \  f_t \ c_{t-1} \ + \ i_t \ z_t & & 
 &\text{cell state} \\
n_t \ &= \  f_t \ n_{t-1} \ + \ i_t  & & 
 &\text{normalizer state} \\
h_t \ &= \ o_t \ \tilde{h}_t \ ,
  &\tilde{h}_t \ &= \ c_t / n_t
&\text{hidden state} \\
z_t \ &= \ \varphi \left( \tilde{z}_t \right) \ , 
  &\tilde{z}_t \ &=  \ \mathbf{w}^\top_{z} \ \mathbf{x}_t \ + \
  r_{z}  h_{t-1} \ + \  b_{z}
  &\text{cell input} \\
i_t \ &= \ \!\exp \left( \tilde{i}_t  \right) \ , 
  &\tilde{i}_t \ &= \ \mathbf{w}^\top_{i} \ \mathbf{x}_t \ + \
  r_{i}  \ h_{t-1} \ + \  b_{i}
  &\text{input gate} \\
f_t \ &= \ \sigma \left(  \tilde{f}_t \right) \ \text{OR} \
\!\exp \left(  \tilde{f}_t \right) \ ,
  &\tilde{f}_t \ &= \ \mathbf{w}^\top_{f} \ \mathbf{x}_t  \ + \
  r_{f}  \ h_{t-1} \ + \  b_{f}
  &\text{forget gate} \\
o_t \ &= \ \sigma \left( \tilde{o}_t \right) \ , 
  &\tilde{o}_t  \ &= \ \mathbf{w}^\top_{o} \ \mathbf{x}_t \ + \
  r_{o}  \ h_{t-1} \ + \  b_{o}
  &\text{output gate} 
\end{align}

We transfer the original LSTM gating techniques, i.e., input- and/or hidden-dependent gating plus bias term, to the new architectures. Exponential activation functions can lead to large values that cause overflows. Therefore, we stabilize gates with an additional state \mbox{$m_t$~\citep{Milakov:18arxiv}}:

\begin{align}
\label{eq:slstmstabil}
m_t \ &= \ \max \left( \log ( f_t ) + m_{t-1} , \log ( i_t ) \right) &\text{stabilizer state} \\
i'_t \ &= \ \!\exp \left( \log \left ( i_t \right) - m_t \right) \ = \ \!\exp \left( \tilde{i}_t - m_t \right) \ \, 
  &\text{stabil. input gate} \\
f'_t \ &= \ \!\exp \left( \log \left( f_t \right) + m_{t-1} - m_t \right) \ \, 
  &\text{stabil. forget gate}
\end{align}

As demonstrated in Appendix~\ref{sec:appsLSTM}, substituting $f_t$ with $f'_t$ and $i_t$ with $i'_t$ during the forward computation leaves both the network’s overall output and the gradients of the loss relative to its parameters unchanged.

\paragraph{Revised Memory Mixing.} Similar to the classical LSTM, the sLSTM architecture allows for several memory cells (refer to Appendix~\ref{sec:appsLSTM}). The presence of multiple memory cells facilitates the mixing of memory through the recurrent links $\mathbf{R}_{\mathbf{z}}$, $\mathbf{R}_{\mathbf{i}}$, $\mathbf{R}_{\mathbf{f}}$, and $\mathbf{R}_{\mathbf{o}}$, which connect the hidden state $\mathbf{h}$ to the input $\mathbf{z}$ of the memory cell and the corresponding gates $\mathbf{i}$, $\mathbf{f}$, and $\mathbf{o}$. Exponential gating introduces a novel characteristic to the process of memory mixing. In this revised sLSTM formulation, it is possible to allocate multiple heads, each capable of memory mixing internally, although such mixing does not occur between different heads. The addition of heads, paired with exponential gating mechanisms in sLSTM, represents a fundamentally new mechanism for combining memory content.

\subsection{mLSTM}
\label{sec:mLSTM}

In order to boost the storage capacity of LSTMs, we generalize the standard memory cell from a single scalar value $c \in \mathbb{R}$ to a full matrix representation $\mathbf{C} \in \mathbb{R}^{d \times d}$. With this modification, memory retrieval relies on matrix multiplication. Specifically, at each time step $t$, the model is tasked with storing a pair of vectors: a key vector $\mathbf{k}_t \in \mathbb{R}^d$ and an associated value vector $\mathbf{v}_t \in \mathbb{R}^d$, following the nomenclature commonly used in Transformers. At a subsequent time $t+\tau$, a query vector $\mathbf{q}_{t+\tau} \in \mathbb{R}^d$ should allow for the recovery of the previously stored value $\mathbf{v}_t$. This paradigm aligns with the framework established in Bidirectional Associative Memories (BAMs) \citep{Kohonen:72,Anderson:72,Nakano:72,Anderson:77}.

The covariance update rule~\citep{Sejnowski:77,Dayan:91} for storing a key-value pair is

\begin{align}
\mathbf{C}_t \ &= \ \mathbf{C}_{t-1} \ + \ \mathbf{v}_{t} \ \mathbf{k}_t^\top \ .
\end{align}

We presume the application of layer normalization prior to mapping inputs onto keys and values, resulting in these components possessing zero mean. The covariance-based update mechanism is deemed optimal~\citep{Dayan:91} for maximizing the separability among retrieved binary vectors, aligning with the objective of achieving the highest signal-to-noise ratio. Enhanced separability becomes attainable when retrieval operations are restricted to pairwise interactions, albeit at the cost of quadratic computational complexity, as encountered in attention mechanisms~\citep{Krotov:16,Krotov:17,Ramsauer:21}. Notably, the covariance update procedure is analogous to the approach of Fast Weight Programmers~\citep{Schmidhuber:92ncfastweights,Schlag:21}. Subsequent research has introduced a fixed decay factor to $\mathbf{C}_{t-1}$ and a fixed learning coefficient to 
$\mathbf{v}_{t} \mathbf{k}_t ^\top$~\citep{Ba:16}.
Guided by these developments, we incorporate the covariance update method into the LSTM architecture, wherein the forget gate implements the decay rate, the input gate reflects the learning rate, and the output gate modulates the scaling of the recalled vector.

For this matrix memory, the normalizer state is the weighted sum of key vectors, where each key vector is weighted by the input gate and all future forget gates. Again, the normalizer state keeps record of the strength of the gates. Since the dot product between query and normalizer state can be close to zero, we use the absolute value of this dot product and lower bound it by a threshold (typically 1.0) as done previously~\citep{Sun:23arxiv}. The mLSTM forward pass is:

\begin{align}
\label{eq:mlstm_recurrent_begin}
\mathbf{C}_t \ &= \  f_t \ \mathbf{C}_{t-1} \ + \ 
  i_t \ \mathbf{v}_t \ \mathbf{k}_t^\top & &  &\text{cell state} \\
\mathbf{n}_t \ &= \  f_t \ \mathbf{n}_{t-1} \ + \ 
  i_t \ \mathbf{k}_t & &  &\text{normalizer state} \\
\mathbf{h}_t  \ &= \ \mathbf{o}_t \ \odot \ \tilde{\mathbf{h}}_t
\ , \qquad \qquad 
\tilde{\mathbf{h}}_t \ = \   \mathbf{C}_t \mathbf{q}_t \ / \ 
  \max \left\{ |\mathbf{n}_t^\top \mathbf{q}_t|, 1 \right\} 
   &&&\text{hidden state} \\
\mathbf{q}_t \ &= \ \mathbf{W}_q \ \mathbf{x}_t \ + \ \mathbf{b}_q  & &  &\text{query input} \\
\mathbf{k}_t \ &= \ \frac{1}{\sqrt{d}} \mathbf{W}_k \ \mathbf{x}_t \ + \ \mathbf{b}_k  & &  &\text{key input} \\
\mathbf{v}_t \ &= \ \mathbf{W}_v \ \mathbf{x}_t \ + \ \mathbf{b}_v  & &  &\text{value input} \\
i_t \ &= \ \!\exp \!  \! \left( \tilde{i}_t  \right) \ , 
 \qquad \qquad \ \ \ \ \,
  \tilde{i}_t \ = \ \mathbf{w}^\top_{i} \ \mathbf{x}_t \ + \  b_{i} \ \ \ 
  &&&\text{input gate} \\
f_t \ &= \ \sigma \!  \left(  \tilde{f}_t \right) \ \text{OR} \ \!\exp \!  \! \left(  \tilde{f}_t \right) , \ \ \: \!
\tilde{f}_t \ = \ \mathbf{w}^\top_{f} \ \mathbf{x}_t  \ + \
  b_{f}
  &&& \text{forget gate} \\
\label{eq:mlstm_recurrent_end}
\mathbf{o}_t \ &= \ \sigma \left( \tilde{\mathbf{o}}_t \right) \ , \qquad \qquad \quad \ \ \ \,
  \tilde{\mathbf{o}}_t  \ = \ \mathbf{W}_{\mathbf{o}} \ \mathbf{x}_t \ + \
  \mathbf{b}_{\mathbf{o}}
  &&&\text{output gate} 
\end{align}

Like the original LSTM, mLSTM is capable of supporting several memory cells. In mLSTM, having multiple heads or multiple cells is effectively the same due to the absence of memory mixing. To ensure the stability of the exponential gating mechanisms in mLSTM, we employ the same stabilization approaches utilized for sLSTM (refer to Equation~\ref{eq:slstmstabil}). Given that mLSTM does not involve memory mixing, its recurrence relation can also be expressed in a parallelized form. Additional information can be found in Appendix~\ref{sec:appmLSTM}.

\subsection{xLSTM Architecture}
\label{sec:xLSTMarch}

\paragraph{xLSTM Blocks.}
\label{sec:xLSTMblock}
An xLSTM block is intended to create a nonlinear high-dimensional representation of previous inputs, thereby enhancing the differentiation among distinct histories or contexts. The ability to discriminate between these histories is essential for accurate prediction of subsequent sequence elements, such as the next token. We build on Cover’s Theorem~\citep{Cover:65}, which implies that patterns, when embedded nonlinearly into higher-dimensional spaces, tend to become more easily separable by linear methods than in their original low-dimensional representations. Two residual block architectures are explored: (i) The residual block utilizing post up-projection—akin to the design in Transformers—first produces a nonlinear summary of the sequence history within its original dimensionality, then projects this output linearly to a higher-dimensional space, applies a nonlinear activation, and subsequently projects it linearly back to the original space. This process is illustrated in the left panel of Figure~\ref{fig:Backbones} and the third column of Figure~\ref{fig:xlstm_sketch}. Additional detail is provided in Figure~\ref{fig:appsLSTM_detailed} in the appendix. (ii) The residual block employing pre up-projection, which, similar to State Space Models, performs an initial linear mapping to a high-dimensional representation,

Performs a non-linear condensation of previous information within a high-dimensional space before employing a linear transformation to revert it to its initial domain. When integrating an sLSTM within an xLSTM block, we apply the post up-projection block. Conversely, when utilizing an mLSTM within xLSTM, we select the pre up-projection block due to the greater memory capacity afforded in the higher-dimensional domain. For further clarification, please consult the leftmost diagram in Figure~\ref{fig:Backbones}, the third column of Figure~\ref{fig:xlstm_sketch}, or Figure~\ref{fig:appsLSTM_detailed} located in the appendix.

\paragraph{xLSTM Architecture. }
\label{sec:xLSTMarchitecure}
The xLSTM model is built by stacking blocks in a residual fashion, following the architectures suggested in~\citep{Srivastava:15,He:16}. Our design leverages pre-LayerNorm~\citep{Ba:16b} residual structures, widely employed in modern Large Language Models. Additional visualization is available in the final column of Figure~\ref{fig:xlstm_sketch}.

\subsection{Memory and Speed Considerations}

In contrast to Transformers, xLSTM networks maintain both linear computational costs and fixed memory requirements irrespective of sequence length. Owing to the compressive characteristic of xLSTM memory, these networks are particularly advantageous for use in industrial contexts and resource-constrained edge devices.

The mLSTM variant utilizes a parameter-free memory, but incurs high computational cost due to the use of a $d \times d$ matrix memory and corresponding $d \times d$ updates. This design opts for increased memory capacity at the expense of added computational demand. Importantly, these operations can be parallelized efficiently on GPU hardware, which means their practical impact on execution time is minimal.

Whereas mLSTM lends itself to parallel execution similarly to FlashAttention~\citep{Dao:22,Dao:23} and GLA~\citep{Yang:23arxiv}, sLSTM lacks this property because its memory updates involve hidden-to-hidden mixing, which inhibits parallelization. To address this, we engineered an optimized CUDA-based implementation that leverages register-level memory optimization on GPUs, resulting in a runtime that is usually no more than twice that of mLSTM.

\section{Related Work}
\label{sec:related}

\paragraph{Linear Attention.} Various approaches have been proposed to address the Transformer's quadratic time complexity with respect to context length, aiming to scale attention mechanisms linearly. The Synthesizer avoids token-to-token dependency altogether by generating synthetic attention weights~\citep{Tay:20arxiv}. The Linformer employs a low-rank matrix to represent self-attention and further achieves a linear approximation~\citep{Wang:20arxiv}. The Linear Transformer achieves a linear complexity by reformulating the attention operation \citep{Katharopoulos:20}. Meanwhile, the Performer leverages a method based on positive orthogonal random features to obtain a linear approximation of the attention softmax~\citep{Choromanski:21}. Additionally, fast long convolutions have been used in place of attention, notably within the Structured Global Convolution (SGConv)~\citep{Li:22} and Hyena Hierarchy~\citep{Poli:23} models.

\paragraph{State Space Models.} State Space Models (SSMs) have recently garnered significant attention due to their linear complexity with respect to sequence length and their strong empirical results when compared to Transformers. Early work in this area includes the introduction of the Structured State Space sequence model (S4)~\citep{Gu:21}. Subsequent notable developments encompass the Diagonal State Space (DSS) model~\citep{Gupta:22}, Gated State Space (GSS) models~\citep{Mehta:22}, the S5 model~\citep{Smith:22}, Bidirectional Gated SSM (BiGS)~\citep{Wang:22}, H3 model~\citep{Fu:23}, and Mamba~\citep{Gu:24arxiv}.

\paragraph{Recurrent Neural Networks.} Recurrent Neural Networks (RNNs) have been proposed as alternatives to Transformer architectures and attention mechanisms, primarily because of their favorable linear scaling with respect to context length. Recent work on Deep Linear Recurrent Units (LRUs) has demonstrated strong performance in language modeling tasks~\citep{Orvieto:23, De:24arxiv}, and similar gains have been observed with the Hierarchically Gated Linear RNN (HGRN)~\citep{Qin:23} as well as its successor, HGRN2~\citep{Qin:24arxiv}. The RWKV model~\citep{Peng:23arxivshort,Peng:24arxivshort} represents a prominent RNN-based methodology for large language modeling, achieving results that are competitive with Transformer-based models.

\paragraph{Gating.}
Gating represents a foundational concept in LSTM architectures, a principle that has been independently revisited and adapted by several contemporary methods. Notable examples where gating mechanisms are central include HGRN~\citep{Qin:23}, HGRN2~\citep{Qin:24arxiv}, Gated Linear Attention (GLA)~\citep{Yang:23arxiv}, Gated State Space (GSS) models~\citep{Mehta:22}, Bidirectional Gated SSM (BiGS)~\citep{Wang:22}, Moving Average Equipped Gated Attention (MEGA)~\citep{Ma:22}, as well as RWKV~\citep{Peng:23arxivshort} and Mamba~\citep{Gu:24arxiv}.

\paragraph{Covariance Update Rule.}  
To improve the memory capabilities, we augmented the mLSTM cell by incorporating a matrix memory governed by a covariance update rule. Several other approaches employ similar update strategies, such as Fast Weight Programmers~\citep{Schmidhuber:92ncfastweights,Schlag:21}, RWKV-5 and RWKV-6~\citep{Peng:24arxivshort}, Retention~\citep{Sun:23arxiv}, Linear Transformer~\citep{Katharopoulos:20}, and HGRN2~\citep{Qin:24arxiv}.

\paragraph{Most Related.} xLSTM bears the strongest resemblance to models such as Retention~\citep{Sun:23arxiv}, RWKV~\citep{Peng:23arxivshort,Peng:24arxivshort}, and HGRN2~\citep{Qin:24arxiv}, which similarly employ matrix memory and/or gating mechanisms. Notably, unlike the recently introduced sLSTM, these models lack support for memory mixing. This capability—unique to sLSTM—enables the model to effectively handle state tracking challenges. As a result, LSTMs exhibit greater expressive power compared to State Space Models (SSMs) and Transformers~\citep{Merrill:24,Deletang:23}. Effective state tracking is essential, for example, in scenarios that require program evaluation or the monitoring of entities across extensive narratives.

\paragraph{Residually Stacking Architectures.} The xLSTM architectures, consistent with the majority of modern, large-scale deep learning models, are built by stacking foundational modules with residual connections~\citep{Srivastava:15,He:16}.

This modular stacking paradigm has paved the way for deep convolutional neural networks~\citep{He:16} and, notably, for Transformer models~\citep{Vaswani:17}. The Transformer architecture underpins the development of Large Language Models (LLMs) such as GPT-3~\citep{Brown:20short}, ChatGPT~\citep{Schulman:22short}, GPT-4~\citep{Achiam:23gpt4short}, Megatron-LM~\citep{Shoeybi:19}, Gopher~\citep{Rae:21short}, ERNIE 3.0 Titan~\citep{Wang:21arxivshort}, GLaM~\citep{Du:21glamshort}, Chinese M6~\citep{Lin:21}, AlexaTM 20B (multilingual)~\citep{Soltan:22}, OPT~\citep{Zhang:22opt}, Chinchilla~\citep{Hoffmann:22short}, BLOOM~\citep{Scao:22arxivshort}, GLM-130B~\citep{Zeng:22arxivshort}, LaMDA~\citep{Thoppilan:22short}, PaLM~\citep{Chowdhery:22short}, Llama~\citep{Touvron:23arxiv}, and Gemini~\citep{GeminiTeam:23arxiv,Reid:24arxiv}.

\section{Experiments}
\label{sec:Exp}

This section presents an empirical assessment of xLSTM, benchmarking its performance in language modeling against established approaches. Section~\ref{sec:ExpSyntethic} examines xLSTM's abilities using synthetic benchmark tasks. In Section~\ref{sec:ExpComparison}, we measure and compare the validation perplexity of several contemporary language modeling techniques, all trained with 15B tokens from the SlimPajama dataset~\citep{Soboleva:23}, alongside conducting ablation analyses of xLSTM. Additionally, we analyze the scaling properties of these approaches following the methodologies outlined by \citet{Kaplan:20} and \citet{Brown:20short}. Subsequently, Section~\ref{sec:ExpLanguage} details an extensive language modeling evaluation, where both xLSTM and the leading contenders from Section~\ref{sec:ExpComparison} are trained with an expanded dataset comprising 300B tokens from SlimPajama~\citep{Soboleva:23}. In this part, we probe each method's ability to generalize to longer sequences, evaluate validation perplexity as well as downstream task performance~\citep{Sutawika:24}, test across 571 domains contained in the PALOMA language benchmark~\citep{Magnusson:23arxivshort}, and revisit scaling behavior analysis utilizing a training set twenty times larger than in previous experiments.

\label{sec:xlstm_blocks_notation}
Throughout all experiments, we denote xLSTM[$a$:$b$] as representing a configuration where the ratio of mLSTM-based blocks to sLSTM-based blocks is $a/b$. For instance, xLSTM[7:1] specifies a total of eight blocks, of which seven employ mLSTM and one utilizes sLSTM. When working with a standardized block count of 48, this ratio yields 42 blocks based on mLSTM and 6 blocks based on sLSTM. In addition, in every experiment, mLSTM and sLSTM blocks are always accompanied by pre and post up-projection blocks, respectively.

\subsection{Synthetic Tasks and Long Range Arena}
\label{sec:ExpSyntethic}
Initially, we evaluate xLSTM's novel exponential gating mechanism combined with memory mixing on formal language benchmarks~\citep{Deletang:23}. Subsequently, the capabilities of xLSTM's updated matrix memory are analyzed using the Multi-Query Associative Recall task~\citep{Arora:23arxiv}. Lastly, we investigate xLSTM's ability to handle lengthy input sequences by measuring its performance on the Long Range Arena~\citep{Tay:21}.

\paragraph{Evaluation of xLSTM’s Exponential Gating with Memory Mixing.} We assess the performance of xLSTM’s recently introduced exponential gating with memory mixing, designed to address state tracking challenges~\citep{Merrill:24, Merrill:23}. Building upon and extending the formal language benchmarks from \citet{Deletang:23}, we adapt them for multi-length training to facilitate evaluation on length extrapolation. Further methodological details and comprehensive results for all tasks are provided in Appendix~\ref{sec:appExpSynthetic-formalLang}. Our benchmarking study compares xLSTM against a variety of alternative approaches, including Transformers, State Space Models, and traditional Recurrent Neural Networks. Performance for each model is quantified by computing the accuracy specifically on those tokens that are essential to the task, with this metric normalized between 0 (chance level) and 1 (exact match). In particular, we evaluate 2-block versions of the following architectures on these formal language tasks: xLSTM[0:1] (utilizing solely the sLSTM block), xLSTM[1:0] (utilizing only the mLSTM block), xLSTM[1:1], Llama, Mamba, RWKV, Retention, Hyena, the vanilla LSTM, and the LSTM utilized within Transformer blocks (LSTM (Block)). The outcomes of these experiments are depicted in Figure~\ref{fig:formal-main}. Approaches such as Transformers and State Space Models, when deployed without memory mixing (and thus lacking state tracking capabilities), are unable to solve certain regular grammar problems, such as the parity task. This finding is consistent with existing research indicating that Transformers and State Space models exhibit inherently weaker computational power than RNNs~\citep{Merrill:24, Merrill:23, Deletang:23}.

\paragraph{Test of xLSTM's Memory Capacities on Associative Recall Tasks.} This study evaluates the matrix-based memory mechanism in xLSTM by examining its performance on the Multi-Query Associative Recall benchmark~\citep{Arora:23arxiv}. In this setting, each input sequence is comprised of randomly assigned key-value pairs from an extensive vocabulary, requiring the model to memorize these associations for subsequent retrieval. To increase the challenge beyond the standard task, both the quantity of key-value pairs—up to 256—and the maximum sequence length—up to 2048 tokens—are raised. These enhancements permit a more comprehensive examination of the memory limitations of competing models. Comparative analysis involves 2-block configurations from Llama, Mamba, RWKV-5, RWKV-6, xLSTM[1:1], and xLSTM[1:0]. Model performance is assessed based on their recall accuracy for the target pairs. Due to their exponential memory capacity relative to the embedding size~\citep{Ramsauer:21}, Transformer-based models like Llama serve as the performance benchmark for this evaluation. Findings are depicted in Figure~\ref{fig:mqar-main}. xLSTM[1:1] demonstrates the highest recall rate among non-Transformer models, even at smaller model sizes. Notably, incorporating the sLSTM block does not impair memory abilities; on the contrary, it enhances recall, particularly at the highest task difficulty level with 256 pairs. Further experimental details and analyses are available in Appendix~\ref{sec:appExpSynthetic-mqar}, where extrapolation studies suggest that the expanded memory of xLSTM supports generalization to contexts surpassing those encountered during model training.

\paragraph{Test of xLSTM's Long Context Capabilities on Long Range Arena.} In order to evaluate how well xLSTM manages extended sequences and substantial context lengths, we conduct a comparison with other techniques using the Long Range Arena~\citep{Tay:21}. Across all examined tasks, xLSTM delivers reliably strong results, indicating that the xLSTM architecture is highly effective at addressing various challenges associated with long context scenarios.

For more details, see Appendix~\ref{sec:appExpSynthetic-lra}.

\subsection{Method Comparison and Ablation Study}
\label{sec:ExpComparison}

In order to investigate the primary inquiry of our study—specifically, the capabilities of our proposed LSTM variants at scale for language modeling—we conduct experiments by training xLSTMs, Transformers, State Space Models, and additional approaches using 15B tokens sourced from SlimPajama, all within an auto-regressive framework. Performance is evaluated on the validation dataset, and we further conduct ablation studies focused on xLSTMs.

\paragraph{Comparing xLSTM to Other Methods.} For the purposes of comparison, we train each model on a dataset containing 15B tokens drawn from SlimPajama~\citep{Soboleva:23}. Model evaluation is based on the perplexity achieved on the validation set. We benchmark the following approaches: our novel xLSTM; GPT-3 (Transformer)~\citep{Brown:20short}; Llama (Transformer)~\citep{Touvron:23arxiv}; H3 (SSM)~\citep{Fu:23}; Mamba (SSM)~\citep{Gu:24arxiv}; RWKV-4, RWKV-5, and RWKV-6 (all RNNs)~\citep{Peng:23arxivshort, Peng:24arxivshort}; GLA (a linear Transformer)~\citep{Yang:23arxiv}; HGRN and HGRN2 (RNNs)~\citep{Qin:23, Qin:24arxiv}; RetNet (linear Transformer)~\citep{Sun:23arxiv}; Hyena (linear Transformer)~\citep{Poli:23}; and both xLSTM[1:0] and xLSTM[7:1], as defined in Section~\ref{sec:xlstm_blocks_notation}. 

Training was performed using mixed precision for all models. Note that for RWKV-5, RWKV-6, GLA, and HGRN2, mixed-precision was not executed via PyTorch's automated mixed-precision feature (additional details are given in Appendix Section~\ref{sec:appGenTrainProc}).

We organize the methods under the following categories: (a) Transformers, (b) State Space Models (SSMs), and (c) Recurrent Neural Networks (RNNs), in addition to linear Transformers. Linear Transformers are a family of architectures that replace traditional attention mechanisms in Transformers with computationally linear operations. All compared models are parameter-matched to a GPT-3 configuration, consisting of 350M parameters with an embedding size of 1024 and 24 residual blocks. Notably, only GPT-3 employs shared parameters between token and output embeddings, resulting in a marginally lower parameter count.

As summarized in Table~\ref{tab:spaj15b_model_results}, xLSTM achieves superior validation perplexity relative to the other evaluated approaches. Appendix~\ref{sec:appExpComparison} provides further discussion and supplementary findings. Scaling trends visualized in Figure~\ref{fig:temp_sclaw15B} further suggest that xLSTM retains its performance edge as model size increases.

\paragraph{Ablation Studies.}
As shown in Table~\ref{tab:spaj15b_model_results} and Figure~\ref{fig:temp_sclaw15B}, xLSTM attains strong performance in language modeling tasks after training on 15B tokens sourced from SlimPajama. To systematically analyze the contribution of each modification transforming LSTM into xLSTM, we incrementally adapt the standard LSTM architecture. The first step involves embedding LSTM layers within pre-LayerNorm residual frameworks. Next, we augment the structure by attaching a post up-projection module. In the final stage, we incorporate exponential gating mechanisms along with matrix memory.

The results are shown in Table~\ref{tab:ablstudies} (top). 

Ablation experiments indicate that the notable gains in performance result from the combined effects of exponential gating and the matrix memory component. Given the significance of gating in both RNNs and State Space Models, various gating approaches are examined through ablation. As shown in Table~\ref{tab:ablstudies} (bottom), making each gate both input-dependent and learnable leads to incremental improvements. Further analysis of the separate backbone elements is provided in Appendix~\ref{sec:appAblDetails}.

\subsection{xLSTM as Large Language Model}
\label{sec:ExpLanguage}

This investigation culminates with comprehensive language modeling experiments designed to evaluate the effectiveness of xLSTM as a large language model (LLM). For these purposes, we expand the training corpus to 300B tokens sourced from SlimPajama, aligning our training scale with that of other contemporary approaches, such as Mamba~\citep{Gu:24arxiv} and Griffin~\citep{De:24arxiv}.

The performance of xLSTM is contrasted with RWKV-4, Llama, and Mamba, each serving as the exemplar for its respective method category described in Section~\ref{sec:ExpComparison}.
RWKV-4 is used as the representative RNN baseline since suitable training precision configurations for RWKV-5, RWKV-6, and HGRN2 were only established after commencement of the 300B token experiments (see Appendix~\ref{sec:appGenTrainProc}).
A variety of model scales are explored (125M, 350M, 760M, 1.3B parameters), examining all models for their ability to extrapolate to longer sequences and determining their accuracy on the validation set. Furthermore, the models are subjected to downstream evaluation, language modeling assessment across 571 domains drawn from the PALOMA benchmark, and analyses of their scaling laws.

\paragraph{Sequence Length Extrapolation.}
\label{sec:spaj300B_ctx_extrapolation}
To begin, we examine how 1.3B-parameter variants of xLSTM, RWKV-4, Llama, and Mamba handle sequence length extrapolation. These models are each trained using a context window of 2048 tokens, after which they are evaluated at longer context lengths extending up to 16384 tokens. The comparative results can be found in Figure~\ref{fig:spaj300B_seq_extrapolation}. Notably, xLSTM stands out from the alternatives by consistently achieving low perplexity scores as the context length increases.

\paragraph{Validation Perplexity and Downstream Tasks.}
\label{sec:spaj_downstream_eval}
Next, we assess xLSTM, RWKV-4, Llama, and Mamba models—across all examined parameter scales—on the SlimPajama validation benchmark for the task of next token prediction. In addition, we consider their effectiveness on downstream tasks oriented towards common sense reasoning abilities.

The third column in Table~\ref{tab:lmeval} presents the validation set perplexities associated with the various methods evaluated. Among these, xLSTM[1:0] and xLSTM[7:1] consistently achieve the lowest perplexities across every model size. The remaining columns in Table~\ref{tab:lmeval} show the results on downstream tasks. Across almost all tasks and model sizes, xLSTM outperforms the alternatives—except for the ARC task, where Mamba sometimes yields the highest performance. Further details are available in Appendix~\ref{sec:appExpLanguage}.

\paragraph{Performance on PALOMA Language Tasks.} Third, across every model scale, we assess the next token prediction accuracy of xLSTM, RWKV-4, Llama, and Mamba on the PALOMA language benchmarks~\citep{Magnusson:23arxivshort}. The evaluation relies on next token perplexity computed over 571 distinct text domains, spanning from sources such as nytimes.com to Reddit forums like r/depression.

Table~\ref{tab:ppleval} presents the perplexity results for these models, segmented into general language modeling domains (first seven columns) and specific fine-grained domains (last five columns). On this suite of tasks, xLSTM[1:0] consistently outperforms xLSTM[7:1].

When compared to other models, xLSTM[1:0] achieves lower perplexity scores than Mamba in 568 of 571 domains (99.5\%), outperforms Llama in 486 domains (85.1\%), and bests RWKV-4 in 570 domains (99.8\%). Further experimental details are provided in Appendix~\ref{sec:appExpLanguage}.

\paragraph{Scaling Laws.} Fourth, we investigate the scaling properties characterized by a power-law relationship, which enables predictions about performance as model size increases~\citep{Kaplan:20,Brown:20short}. Figure~\ref{fig:temp_sclaw300B} depicts this scaling trend. Although all examined models exhibit comparable scaling patterns, their baseline performances differ. RWKV-4 exhibits the lowest performance, succeeded by Llama and then Mamba. The xLSTM outperforms Mamba, maintaining a lead similar in magnitude to that which Mamba holds over Llama. The observed scaling trends suggest that xLSTM is likely to sustain or even improve its relative performance over Transformers and State-Space models as model size grows.

\textbf{Generation Times and Maximal Throughput.} In Figure~\ref{fig:inference_time_generation}, we evaluate the generation speed for our xLSTM variants at the 1.3B parameter scale, while Figure~\ref{fig:inference_time_decoding_throughput} (left) illustrates their peak throughput. Our benchmarking includes similarly sized implementations of Mamba, Llama, and RWKV from HuggingFace, with the Llama comparison featuring a static key-value cache. It should be noted that, during our experimentation, neither full cache compilation for Transformer Models nor compilation of Mamba using \texttt{torch.compile} were operational. All text generation assessments were conducted using a batch size of 1 and a pre-fill of 16 across all models—a setting that is especially advantageous for the Transformer. The results in Figure~\ref{fig:inference_time_generation} highlight that xLSTM, along with the other recurrent architectures (Mamba and RWKV-4), exhibits linear time scaling, whereas Llama demonstrates quadratic complexity. Decoding throughput evaluations account for a variety of batch sizes and pre-fill configurations for the Llama model. Figure~\ref{fig:inference_time_decoding_throughput} (right) indicates that, because xLSTM maintains constant memory usage, it accommodates significantly larger batch sizes compared to Llama and consequently attains superior throughput.

\section{Limitations}

(i) Unlike mLSTM, the architecture of sLSTM prevents parallelizable memory mixing, making it unsuitable for highly parallel execution. However, we have implemented a CUDA kernel specifically for sLSTM that operates at less than twice the runtime of our parallel mLSTM version. (ii) Our mLSTM CUDA kernels have not undergone significant optimization, leading to implementations that are approximately four times slower than either FlashAttention or Mamba's scan. By drawing on techniques similar to those used in FlashAttention, improved, more efficient CUDA kernels could be developed. (iii) The use of matrix memory in mLSTM introduces substantial computational demands because matrices of size $d \times d$ are involved. Nevertheless, since both memory reading and writing steps are parameter-free and employ parallelizable standard matrix computations, the added time cost for memory complexity remains modest. (iv) Appropriate initialization of the forget gates is a critical consideration. (v) Because the matrix memory’s capacity does not scale with the sequence length, larger sequence lengths risk saturating memory when dealing with extended contexts. Despite this, for contexts of up to 16k tokens, this has not shown to pose a practical concern (refer to Section~\ref{sec:spaj300B_ctx_extrapolation}). (vi) Given the significant computational demands of large-scale language experiments, we have not yet engaged in comprehensive architectural or hyperparameter optimization—particularly for the larger xLSTM models. Unlocking the complete capabilities of xLSTM will likely necessitate a dedicated and extensive tuning process.

\section{Conclusion}
We have provided a partial response to our central inquiry: To what extent can language modeling be advanced by scaling LSTM architectures to the order of billions of parameters? Our findings suggest that the answer is: ``At a minimum, to the same extent as contemporary approaches such as Transformers and State Space Models.'' In this work, we introduced the xLSTM, extending the traditional LSTM framework through exponential gating mechanisms combined with memory mixing and a reimagined memory architecture. Our experiments demonstrate that xLSTM models yield competitive results in language modeling benchmarks when compared to leading frameworks like Transformers and State Space Models. The observed scaling laws imply that scaling xLSTM models to even larger sizes could present them as strong alternatives to prevalent Large Language Models based on the Transformer paradigm. Beyond language modeling, xLSTM holds promise for advancing various domains, including Reinforcement Learning, Time Series Prediction, and physical systems modeling.

\end{document}